\section{Why bundle proofs with data at all?}\label{sec:06-groups:06-why-bundle:1}

Given a term \texttt{Grp : Group G}, \emph{any} theorem proved about a ``generic'' \texttt{Group G} (using only \texttt{Grp.assoc}, \texttt{Grp.id\textbackslash{}\_left}, and so on) automatically applies to \texttt{intGroup}, to \texttt{perm3Group} (the previous section's permutation group), and to every other group constructed later (path algebras' underlying additive group, and beyond). This is the payoff of the whole exercise: prove it once, generically, and obtain it for free everywhere. Chapter 7 demonstrates exactly this, applying a generic theorem to a concrete group once it is proved.

\textbf{Mathlib equivalent.} This ``prove it once, get it for free everywhere'' promise is not something Mathlib puts off to a later chapter; it is the reason Mathlib's algebra hierarchy is organized around typeclasses at all. The \emph{same} lemma name applies unchanged to two completely different groups:

\begin{lstlisting}
example (a b c : Int) : (a + b) + c = a + (b + c) := add_assoc a b c
example (f g h : Equiv.Perm (Fin 3)) : (f * g) * h = f * (g * h) := mul_assoc f g h
\end{lstlisting}

\href{https://loogle.lean-lang.org/?q=add_assoc}{\texttt{add\textbackslash{}\_assoc}}/\href{https://loogle.lean-lang.org/?q=mul_assoc}{\texttt{mul\textbackslash{}\_assoc}} were proved exactly once, generically over \texttt{[AddCommGroup G]}/\texttt{[Group G]}, and both \texttt{Int} and \texttt{Equiv.Perm (Fin 3)} obtain the fact automatically simply by having a \texttt{Group}/\texttt{AddCommGroup} instance. Nothing about \texttt{Int} or permutations is re-proved at either call site. This is the library-scale version of the payoff Chapter 7 walks through by hand for \texttt{perm3Group}.
